% ============================================================
%  4. METHODOLOGY
% ============================================================
\chapter{Methodology}\label{sec:methodology}



\section{Forecasting Framework}

All models are estimated on a training sample running from December 2000 to March 2024 and evaluated on a hold-out test set covering the final 24 months from April 2024 to March 2026. Model selection and specification are conducted exclusively on the training sample to avoid look-ahead bias. As a robustness check, Time Series Cross-Validation (TSCV) is applied using an expanding window with a minimum training length of 60 months, averaging one-step-ahead forecast errors across the full sample.


\section{ARIMA Benchmark}

The univariate ARIMA model serves as the benchmark, forecasting Danish inflation using only its own past values and imposing no economic structure. The general seasonal ARIMA$(p,d,q)(P,D,Q)_{12}$ model is written as:
\[
\Phi_P(B^{12})\,\phi_p(B)\,\Delta^d\Delta^D_{12}\,y_t 
= \Theta_Q(B^{12})\,\theta_q(B)\,\varepsilon_t
\]
where $B$ is the backshift operator, $\phi_p$ and $\theta_q$ are the non-seasonal AR and MA polynomials, $\Phi_P$ and $\Theta_Q$ are their seasonal counterparts, and $\varepsilon_t$ is white noise. Order selection follows \textcite{box_jenkins_1976} and is automated by minimising AICc over a grid of candidate specifications \parencite{hyndman_athanasopoulos_2021}. Residual adequacy is assessed using the Ljung-Box test.

% \section{Structural Breaks}

% The QLR test identified two structural breaks in Danish inflation: December 2012, coinciding with the tail end of the European sovereign debt crisis, and August 2021, marking the onset of the post-COVID inflation surge. Both breaks are incorporated into all models as step dummy variables, taking the value zero before the break date and one thereafter. Omitting these dummies would cause the models to absorb permanent level shifts through the error structure, distorting parameter estimates and degrading forecast performance. \textcolor{red}{er det her dobbelt confekt???}



% \section{Dynamic Regression with ARIMA Errors}

% The dynamic regression model extends the ARIMA benchmark by incorporating external predictors. Danish inflation is modelled as a function of lagged first-differenced unemployment and EU27 inflation, with the remaining serial dependence in the errors captured by an ARIMA process:

% \[
% y_t = \mathbf{x}_t'\boldsymbol{\beta} + \eta_t, \qquad 
% \phi_p(B)\,\eta_t = \theta_q(B)\,\varepsilon_t
% \]

% where $y_t$ is Danish inflation, $\mathbf{x}_t$ contains the external regressors and structural break dummies, and $\eta_t$ follows an ARIMA$(p,0,q)$ process. This formulation tests whether unemployment and European inflation contain predictive information for Danish inflation beyond its own autocorrelation structure.



\section{Dynamic Regression with ARIMA Errors}

The dynamic regression model extends the ARIMA benchmark by incorporating external predictors while preserving the model's ability to capture serial dependence. The general specification is:
\[
y_t = \mathbf{x}_t'\boldsymbol{\beta} + \eta_t, \qquad
\phi_p(B)\,\eta_t = \theta_q(B)\,\varepsilon_t,
\]
where $y_t$ is Danish inflation, $\mathbf{x}_t$ contains the external regressors and structural break dummies, and $\eta_t$ follows an ARIMA$(p,0,q)$ process selected by AIC minimisation. 

In this paper we will extend our benchmark ARIMA model with three different models, each testing a specific hypothesis about improving our forecast of the Danish inflation.

\textbf{Model A} tests the domestic Phillips Curve — the hypothesis that lower unemployment drives higher inflation — by regressing Danish inflation on first-differenced Danish unemployment at lags 0 to 3.
\[
y_t = \beta_0 \cdot d_{2013} + \beta_1 \cdot d_{2021} +
\sum_{j=0}^{3} \gamma_j\,\Delta u^{\text{DK}}_{t-j} + \eta_t.
\]
Lags 0 to 3 are included to allow for delayed transmission of labour market conditions to prices, capturing effects up to one quarter ahead.

\textbf{Model B} asks whether the domestic Phillips Curve is sufficient for a small open economy. Denmark's ERM~II peg ties monetary conditions closely to the euro area, suggesting that European price pressures may spill over into Danish inflation independently of the domestic labour market. 

Model B therefore extends Model A by additionally including EU27 inflation at lags 0--3, testing whether the open-economy channel improves out-of-sample forecast accuracy:
\[
y_t = \beta_0 \cdot d_{2013} + \beta_1 \cdot d_{2021} +
\sum_{j=0}^{3} \gamma_j\,\Delta u^{\text{DK}}_{t-j} +
\sum_{j=0}^{3} \delta_j\,\pi^{\text{EU27}}_{t-j} + \eta_t.
\]
\textbf{Model C} addresses the risk of overfitting inherent in Models A and B, where lag structures are pre-specified rather than data-driven. Variable and lag selection instead follows the General-to-Specific (GETS) methodology \parencite{doornik2014}: a General Unrestricted Model (GUM)is specified with up to six lags of first-differenced Danish unemployment, EU27 inflation, and first-differenced EU27 unemployment, together with the two structural break dummies. The GUM is then reduced by sequentially eliminating statistically insignificant regressors at the five percent level. Impulse Indicator Saturation (IIS) and Step Indicator Saturation (SIS) are applied simultaneously to detect individual outliers and permanent level shifts beyond the pre-specified break dummies. The retained specification is carried forward as Model C. 





\section{Forecast Evaluation}

Forecast accuracy is evaluated using the Root Mean Squared Error (RMSE) and Mean Absolute Error (MAE), defined as:
\[
\text{RMSE} = \sqrt{\frac{1}{h}\sum_{t=1}^{h}(y_t - \hat{y}_t)^2}, \qquad
\text{MAE}  = \frac{1}{h}\sum_{t=1}^{h}|y_t - \hat{y}_t|
\]
where $h = 24$ is the forecast horizon and $\hat{y}_t$ is the model forecast. RMSE penalises large errors more heavily due to the quadratic term, while MAE treats all errors equally. A model is considered to improve upon the ARIMA benchmark if it achieves a lower RMSE and MAE over the test period.